\chapter{Numerical evidence and applications}\label{ch:evidence}
The numerical record has four distinct roles: exact-law controls test values with independent formulas, heterogeneous spectra test inversion across irregular coefficients, simulated kernels test the operator pipeline, and timing measurements characterize particular implementations. This chapter keeps those roles separate. The source package records the inputs and environments behind the retained experiments. Current regression and containment checks are reported in Appendix~\ref{app:reproduction}.

\section{Deep-tail controls}\label{sec:final-3}
A directed enclosure includes rounding as well as analytic remainders. An analytic ledger evaluated in ordinary floating point does not. The labels in this chapter identify which conclusion a result supports. The retained timing measurements are not presented as a new hardware comparison.

The method is designed for a regime in which direct survival probabilities can
underflow long before the mathematical probability becomes numerically
unresolvable.  The engine therefore treats \(\log p\) as the primary numerical
quantity and keeps approximation certificates alive in the same extreme-tail
regime.

The validation ladder spans seven qualitatively different settings:

\begin{center}
\small
\begin{tabular}{@{}P{0.27\linewidth}P{0.13\linewidth}P{0.22\linewidth}P{0.16\linewidth}P{0.12\linewidth}@{}}
\toprule
validation regime & deepest $\log_{10}p$ & observed error & certificate & spectral size\\
\midrule
directed-rounding $0.02\chi^2_{80}$ control
& $-5000$ & enclosure width $3.34\times10^{-11}$ decades
& multiplicative ratio $\le 1+7.70\times10^{-11}$ & 80 df\\
large-scale positive spectrum
& $-922.031162$ & below stored reference resolution & $5.71\times10^{-9}$ & rank 1997\\
strong-dependence kernel
& $-457.973916$ & below stored reference resolution & $5.20\times10^{-9}$ & rank 996\\
heterogeneous inversion benchmark near $10^{-16}$
& $-16$ & $2.88$--$6.52\times10^{-15}$ decades & independent integral converged & ranks 80--120\\
signed finite-endpoint exact control
& $-200$ & $\le 2.84\times10^{-14}$ decades & $7.23\times10^{-13}$ & rank 80\\
signed mixed multiscale
& $-81.891937$ & $1.42\times10^{-12}$ decades vs full grid & $5.78\times10^{-9}$ & rank 340\\
noncentral exact controls
& $-40$ & $\le 9.26\times10^{-15}$ decades & passing & up to 50 df\\
\bottomrule
\end{tabular}
\end{center}

The analytic scaled-chi-square ladder also includes
\[
\log_{10}p=-100,-300,-1000,-5000,
\]
with respective reported absolute log-probability discrepancies
\[
2.84\times10^{-14},\quad
<10^{-15},\quad
3.41\times10^{-13},\quad
<10^{-15}
\]
in the stored benchmark output.  The corresponding analytic relative certificates
remain finite throughout, reaching only
\[
3.85\times10^{-11}
\]
at \(10^{-5000}\).

A large-scale run with \(n=2000\), \(m=50{,}000\), 49,951 retained
features, and quadratic-form rank 1997 produced
\[
\boxed{\log_{10}p=-922.0311624669905}
\]
using 79 spectral packets and \(K=128\).  Its independently evaluated
full-spectrum tilted reference agreed within the stored reporting precision, while the
reported relative certificate was
\[
5.71\times10^{-9}.
\]

The corrected strong-dependence case study independently reached
\[
\boxed{\log_{10}p=-457.9739160722001}
\]
with a finite \(5.20\times10^{-9}\) certificate.  The signed finite-endpoint
control passed its analytic ledger through \(10^{-200}\), while a rank-340 mixed-sign
multiscale spectrum passed its analytic ledger to
\[
\boxed{\log_{10}p=-81.89193686056221}.
\]
Finally, the noncentral engine matched exact noncentral-chi-square controls
from \(10^{-4}\) through \(10^{-40}\), with maximum observed discrepancy
\[
9.26\times10^{-15}
\]
decades.

\begin{figure}[htbp]
\centering
\includegraphics[width=.86\linewidth]{figure_23_deep_tail_precision_envelope.png}
\caption{Demonstrated precision envelope across positive analytic controls,
a large-scale rank-1997 spectrum, corrected strong-dependence spectra,
signed finite-endpoint and mixed-sign forms, and noncentral exact controls.  Points whose stored benchmark difference rounded to zero are
placed at $10^{-16}$ decades only to make them visible on the logarithmic
vertical axis.}
\label{fig:deep-tail-precision-envelope}
\end{figure}

\subsection{Independent characteristic-function inversion near \texorpdfstring{\(10^{-16}\)}{1e-16}}

Closed-form controls are important, but they can hide numerical difficulties
that appear only for heterogeneous spectra.  An independent inversion test was
therefore run on four generalized chi-square laws near
\[
P(Q\ge q)\approx10^{-16}.
\]

The reference evaluates the original characteristic function directly through
the central Imhof inversion identity
\[
P(Q\ge q)
=
\frac12
+
\frac1\pi
\int_0^\infty
\frac{\sin\Theta(u)}{u\,R(u)}\,du,
\]
where, for distinct positive coefficients \(\lambda_j\) with multiplicities
\(h_j\),
\[
\Theta(u)
=
\frac12\sum_j h_j\arctan(\lambda_j u)
-
\frac12qu,
\qquad
R(u)
=
\prod_j
\left(1+\lambda_j^2u^2\right)^{h_j/4}.
\]
This route does not use exponential tilting, the COS expansion, spectral
packets, or the K-COS error controller.

The four spectra include balanced, skewed, dominant-spike, and higher-rank
cases.  Each reference integral was evaluated at two truncation limits.  The
omitted integral after \(U\) was bounded using
\[
\frac1\pi
\int_U^\infty\frac{du}{uR(u)}
\le
\frac{C}{\pi p}U^{-p},
\qquad
p=\frac12\sum_jh_j,
\qquad
C=\prod_j\lambda_j^{-h_j/2}.
\]

\begin{center}
\small
\begin{tabular}{@{}lrrrr@{}}
\toprule
spectrum & rank & \(|\Delta\log_{10}p|\) & relative difference & reference tail bound / \(p\)\\
\midrule
balanced & 80 & $6.52\times10^{-15}$ & $1.50\times10^{-14}$ & $6.84\times10^{-20}$\\
skewed & 80 & $4.13\times10^{-15}$ & $9.50\times10^{-15}$ & $4.44\times10^{-18}$\\
spiky & 80 & $5.19\times10^{-15}$ & $1.20\times10^{-14}$ & $7.00\times10^{-17}$\\
balanced & 120 & $2.88\times10^{-15}$ & $6.64\times10^{-15}$ & $4.19\times10^{-29}$\\
\bottomrule
\end{tabular}
\end{center}

The two inversion truncations agree to between approximately
\[
1.2\times10^{-21}
\quad\text{and}\quad
3.7\times10^{-17}
\]
decades across the four cases.  Thus the slow reference is numerically stable
at a scale well below the observed K-COS-versus-inversion difference.
Agreement between truncations is convergence evidence. The omitted-integral
bound alone does not enclose quadrature discretization or numerical rounding,
so this comparison is not assigned the directed-enclosure guarantee.

\begin{figure}[htbp]
\centering
\includegraphics[width=.82\linewidth]{figure_27_imhof_p1e16_validation.png}
\caption{Independent direct-inversion validation near $p=10^{-16}$ on four
heterogeneous spectra.  Circles show the relative K-COS/inversion difference. Crosses show the analytic bound on the inversion integral omitted beyond the
final truncation point.}
\label{fig:imhof-p1e16}
\end{figure}

This test is complementary to the exact scaled-chi-square controls.  It asks
whether the solver remains accurate when no single-gamma reduction is
available and when an algorithmically different inversion of the full
characteristic function is used as the reference.

These results have different evidentiary roles.  The scaled-chi-square,
finite-endpoint signed, Laplace signed, and noncentral cases provide analytic
controls.  The heterogeneous inversion battery, large-scale, strong-dependence, and
rank-340 mixed-sign cases test irregular spectra for which no single-gamma
closed form is available.  Together they test correctness against known distributions,
representation compression, and stability on the intended high-rank numerical
regime.  The low-rank Laplace example also records a deliberate certificate
failure boundary rather than conflating numerical agreement with certification.


\section{Signed and noncentral controls}
\paragraph{Exact finite-endpoint control.}
For
\[
Q=-0.02\,\chi^2_{80},
\]
the right tail is a lower chi-square tail with finite upper endpoint zero.  The
signed engine agrees with the exact distribution through
\[
\boxed{p=10^{-200}},
\]
with maximum observed log-probability error about $2.84\times10^{-14}$ decades
and certificates around $10^{-12}$.

\paragraph{Mixed-sign high-rank analytic-ledger control.}
A deterministic rank-340 multiscale spectrum with 180 positive and 160
negative coefficients was tested from approximately $10^{-14}$ to
\[
\boxed{p\approx10^{-81.891937}}.
\]
At the deepest point the packetized solver used 35 packets, $K=4096$, and
reported a relative certificate
\[
5.78\times10^{-9}.
\]
Its log-probability differed from the unpacketized full-spectrum signed solver
by about $1.42\times10^{-12}$ decades.  Independent exponential-tilt
importance sampling placed the discrepancy at about $1.41$ Monte Carlo
standard errors.

\paragraph{Retained low-rank limitation.}
The exact mixed-sign identity
\[
Q=0.07(\chi^2_2-\chi^2_2)
\]
is Laplace distributed with scale $0.14$, so for $q\ge0$,
\[
P(Q\ge q)=\frac12e^{-q/0.14}.
\]
The computed values agree with this exact law through $10^{-200}$ to roughly
$10^{-9}$ decades or better.  However, the implemented analytic COS-series
certificate rejects these low-rank mixed cases because the tilted
characteristic function decays only as $|t|^{-2}$.  The implementation reports
this as a certificate failure rather than silently accepting the numerical
value.

\begin{figure}[htbp]
\centering
\includegraphics[width=.84\linewidth]{figure_24_signed_certificate_boundary.png}
\caption{Signed-form certification boundary.  Negative finite-endpoint forms
and the rank-340 mixed spectrum remain certifiable in deep tails.  The exact
rank-4 Laplace control remains numerically accurate but fails the implemented
analytic COS-series certificate because its Fourier decay is much slower.}
\label{fig:signed-certificate-boundary}
\end{figure}


\begin{figure}[htbp]
\centering
\includegraphics[width=.82\linewidth]{figure_22_noncentral_controls.png}
\caption{Independent noncentral controls.  Three scaled isotropic Gaussian
quadratic forms with different dimensions and noncentralities were compared
against SciPy's noncentral-\(\chi^2\) law from \(10^{-4}\) through
\(10^{-40}\).  The largest observed discrepancy was approximately
\(9.3\times10^{-15}\) decades, with all analytic truncation certificates
passing.}
\label{fig:noncentral-controls}
\end{figure}


\section{Rare-variant kernel case study}\label{sec:final-24}

The preceding sections establish operator identities and certified numerical
bounds.  This section tests the same machinery on the project's existing
rare-variant SKAT simulation layer, including explicit LD structure,
covariates, per-variant standardization, rare-variant weighting, and the
phenotype-space score statistic.

\subsection{Simulation provenance and implementation correction}

Two established genotype generators were used.

The first is a latent-Gaussian block-LD generator.  Within a genomic block,
alleles are obtained by thresholding correlated Gaussian latent variables at
variant-specific minor-allele frequencies.

The second is a copy-LD generator in which haplotype states persist along a
block with high probability and are occasionally redrawn. The simulator
must initialize its genotype accumulator to zero before adding haplotypes.
Allocating the accumulator with \verb|np.empty| and then applying
\verb|+=| would make the output depend on unspecified memory contents.
All copy-LD results reported here use explicit zero initialization. Results
from an uninitialized accumulator are excluded from the evidence record.

\subsection{Tested LD regimes}

Four deterministic scenarios were evaluated.

\begin{table}[htbp]
\centering
\caption{Realistic rare-variant simulation scenarios.  Adjacent $r^2$ is
computed within the simulated LD blocks after variant standardization.}
\begin{tabular}{lrrrrr}
\toprule
scenario & $n$ & $m$ & rank & median adjacent $r^2$ & effective rank\\
\midrule
latent LD, moderate & 600 & 3000 & 596 & 0.0069 & 441.9\\
latent LD, strong & 600 & 3000 & 596 & 0.1140 & 240.3\\
copy LD, strong & 600 & 3000 & 596 & 0.8298 & 200.1\\
copy LD, large & 1000 & 5000 & 996 & 0.8179 & 353.0\\
\bottomrule
\end{tabular}
\end{table}

The corresponding spectra therefore span a substantial range of concentration:
stronger LD reduces effective rank even when the nominal quadratic-form rank
remains close to the sample residual dimension.

\begin{figure}[htbp]
\centering
\includegraphics[width=\linewidth]{figures_edition/figure_18_realistic_ld_effective_rank.pdf}
\caption{LD structure changes the effective spectral rank.  The corrected
copy-LD models create much stronger local dependence than the latent-Gaussian
moderate-LD case, while all remain within the same streamed SKAT pipeline.}
\label{fig:realistic-ld-rank}
\end{figure}

\subsection{Streaming directly from raw genotype blocks}

The raw genotype matrix was streamed in resident variant blocks of widths
31, 127, and 509.  Within each block the implementation performs

\[
G_b
\to
G_b-\bar G_b
\to
\text{variance scaling}
\to
M_X G_b
\to
W_b^{1/2}
\to
A_bA_b^\top,
\]

without materializing the full weighted factor.

Across all four LD scenarios and all three resident block sizes, the largest
relative discrepancy against the fully materialized kernel was

\[
\boxed{6.61\times10^{-16}}.
\]

\begin{figure}[htbp]
\centering
\includegraphics[width=\linewidth]{figures_edition/figure_19_realistic_streaming_invariance.pdf}
\caption{Raw-genotype streaming reproduces the materialized SKAT kernel across
multiple LD regimes and block sizes.  The observed discrepancies are at
floating-point summation scale.}
\label{fig:realistic-streaming}
\end{figure}

This test is stronger than applying the streamed identity to a precomputed
factor because the centering, variance scaling, covariate projection, and
rare-variant weights are all performed inside the block loop.

\subsection{Phenotype-space equivalence}

For each scenario, 100 Gaussian null residual draws were evaluated in three
ways. Under the standardized Gaussian null, write $e=M_Xz$ with
$z\sim N(0,I)$. Because $M_XA=A$, the statistic used here is
$Q=\frac12 z^\top Kz$, so its nonzero mixture coefficients are
$\frac12\lambda_j(K)$. The factor $\frac12$ belongs to this case-study
normalization, not to the general quadratic-form definition.

\begin{enumerate}
\item directly through the factor,
\[
Q=\frac12\|A^\top e\|_2^2.
\]
\item through the sample kernel,
\[
Q=\frac12 e^\top Ke.
\]
\item through the eigencoordinates of $K$.
\end{enumerate}

The worst relative discrepancy across all draws and all scenarios was

\[
\boxed{8.32\times10^{-16}}.
\]

Thus the simulated phenotype path, the streamed kernel, and the generalized
$\chi^2$ spectral representation agree to floating-point precision.

\subsection{Deep tails on realistic LD spectra}

For every scenario, tail queries were generated at tilt fractions
0.90, 0.97, and 0.99 of the positive singularity.  K-COS was compared with a
higher-resolution full-spectrum tilted reference.

The large corrected copy-LD case was then pushed further, to fractions
0.995, 0.997, 0.999, and 0.9995.  The deepest tested reference value was

\[
\boxed{\log_{10}p=-457.973916}.
\]

At that point the displayed K-COS and high-resolution full-spectrum values
were identical.  Across the complete realistic-LD battery, the largest
absolute discrepancy was

\[
\boxed{3.21\times10^{-12}\ \text{decades}}.
\]

\begin{figure}[htbp]
\centering
\includegraphics[width=\linewidth]{figures_edition/figure_20_realistic_deep_tail_accuracy.pdf}
\caption{K-COS versus the high-resolution full-spectrum tilted reference on
realistic LD spectra.  Accuracy remains at roughly machine-scale in log
probability through the deepest corrected copy-LD query, near $10^{-458}$.}
\label{fig:realistic-deep-tail}
\end{figure}

For the $n=1000$, $m=5000$ corrected copy-LD spectrum, selected results are

\begin{center}
\begin{tabular}{rrrr}
\toprule
tilt fraction & $\log_{10}p$ & K-COS $K$ & absolute error (decades)\\
\midrule
0.9950 & -65.950669 & 2048 & $4.69\times10^{-13}$\\
0.9970 & -95.305034 & 2048 & $3.70\times10^{-14}$\\
0.9990 & -240.602047 & 8192 & $4.94\times10^{-14}$\\
0.9995 & -457.973916 & 16384 & below stored reference resolution\\
\bottomrule
\end{tabular}
\end{center}

\subsection{Interpretation and boundary}

This simulation block does not add a new algebraic theorem.  It supplies a
more demanding validation layer for the theorem chain already proved.
Specifically, it shows that the exact streaming identities and the K-COS
deep-tail machinery remain numerically stable under substantially different
rare-variant LD structures and through the complete raw-genotype preprocessing
path.

The evidence remains synthetic.  These generators do not reproduce every
population-genetic feature of a sequencing cohort, and this block is not an
external-cohort replication.  It also addresses the Gaussian linear-model null;
finite-sample Bernoulli or logistic calibration remains a separate statistical
problem.



\section{Spectral-packet experiments}
These tables and figures report the retained packet-planner experiments.
Their analytic error bounds exclude floating-point roundoff. They measure a
different numerical path from the directed positive backend, and are not
reinterpreted as directed probability enclosures or new performance runs.

\subsection{Exactly solvable chi-square stress test}

To isolate the effect of exponential tilting, we first used
\[
Q=0.02\chi^2_{80},
\]
whose survival probabilities are available directly from the chi-square distribution. Table \ref{tab:chi} compares a global COS calculation with query-centered tilted COS.

\begin{table}[htbp]
\centering
\caption{Relative error in an exactly solvable chi-square family.}
\label{tab:chi}
\begin{tabular}{rrrr}
\toprule
True tail $p$ & $N$ & Global COS rel. error & Tilted COS rel. error\\
\midrule
$10^{-6}$  & 64  & $3.56\times10^{-8}$ & $2.71\times10^{-14}$\\
$10^{-10}$ & 64  & $1.64\times10^{-4}$ & $6.59\times10^{-15}$\\
$10^{-10}$ & 128 & $2.03\times10^{-6}$ & $9.05\times10^{-15}$\\
$10^{-14}$ & 64  & $1.43$              & $3.53\times10^{-14}$\\
$10^{-14}$ & 128 & $6.15\times10^{-2}$ & $3.87\times10^{-14}$\\
\bottomrule
\end{tabular}
\end{table}

The test demonstrates the numerical motivation for tilting: evaluating a small tail directly under the shifted law is substantially more stable than reconstructing a global CDF and subtracting from one. This observation is not a novelty claim for exponential tilting itself.

\subsection{Genotype-derived spectra}

The second experiment used the existing Komorebi SKAT benchmark generator. Synthetic genotype matrices were centered and scaled, residualized against three covariates, weighted by $(1-\mathrm{MAF})^{24}$, and converted to the positive quadratic-form spectrum
\[
\lambda=\tfrac12\operatorname{eig}(AA^{\mathsf T}).
\]
Three independent $n=400$, $m=1200$ genotype sets yielded numerical rank $397$. Four upper-tail queries per spectrum were generated by choosing saddlepoint locations at fractions $0.70,0.80,0.85,0.90$ of the transform singularity. The comparison reference used exact-spectrum tilted COS with $N=2048$ and a large support parameter $L=20$. The adaptive method was asked for a relative analytic tolerance of $10^{-8}$.

Across the twelve queries:
\begin{itemize}
\item every requested $10^{-8}$ analytic certificate passed.
\item the worst observed relative error was $6.148608\times10^{-11}$.
\item the largest returned analytic relative bound was $7.917203\times10^{-9}$.
\item the planner selected $N\in\{64,128\}$, $L\in\{10,11,12,13\}$, and 11--16 packets.
\end{itemize}

Figure \ref{fig:coverage} plots the observed errors against the returned analytic bounds. Every point lies above the equality line because the vertical coordinate is the bound.

\begin{figure}[htbp]
\centering
\includegraphics[width=0.78\linewidth]{figures_packets/certificate_coverage.png}
\caption{Observed relative error versus the analytic certificate on twelve genotype-derived tail queries. The dashed line is equality.}
\label{fig:coverage}
\end{figure}

\subsection{Tolerance ladder}

At a reference tail probability near $8.99\times10^{-15}$, the same engine was requested to meet three tolerances. Table \ref{tab:ladder} shows how the representation refined automatically.

\begin{table}[htbp]
\centering
\caption{Adaptive refinement at an extreme-tail query.}
\label{tab:ladder}
\begin{tabular}{rrrrrr}
\toprule
Requested $\tau$ & $N$ & $L$ & Packets & Analytic bound & Observed error\\
\midrule
$10^{-6}$  & 128 & 10 & 10 & $3.56\times10^{-7}$  & $3.63\times10^{-12}$\\
$10^{-8}$  & 128 & 12 & 15 & $4.87\times10^{-9}$  & $9.14\times10^{-13}$\\
$10^{-10}$ & 256 & 14 & 24 & $8.21\times10^{-11}$ & $5.65\times10^{-14}$\\
\bottomrule
\end{tabular}
\end{table}

\begin{figure}[htbp]
\centering
\includegraphics[width=0.78\linewidth]{figures_packets/tolerance_ladder.png}
\caption{Requested tolerance, returned certificate, and observed error for one extreme-tail query.}
\label{fig:ladder}
\end{figure}

\subsection{Prototype runtime scaling}

Table \ref{tab:runtime} reports median CPU timings from the Python prototype. ``Adaptive total'' includes planning and analytic error checks. ``Compiled evaluation'' measures reuse of an already compiled spectral passport at the selected grid.

\begin{table}[htbp]
\centering
\caption{Prototype runtime scaling on genotype-derived spectra.}
\label{tab:runtime}
\begin{tabular}{rrrrrr}
\toprule
Rank & Packets & Exact tilted & Adaptive total & Compiled eval. & Compiled speedup\\
\midrule
397  & 13 & 6.66 ms & 10.45 ms & 0.549 ms & $12.1\times$\\
797  & 16 & 13.32 ms & 11.49 ms & 0.604 ms & $22.0\times$\\
1197 & 17 & 21.52 ms & 11.24 ms & 0.651 ms & $33.1\times$\\
\bottomrule
\end{tabular}
\end{table}

\begin{figure}[htbp]
\centering
\includegraphics[width=0.78\linewidth]{figures_packets/runtime_scaling.png}
\caption{Median runtime of exact tilted COS, adaptive K-COS with analytic ledger including planning, and evaluation of an already compiled passport.}
\label{fig:runtime}
\end{figure}

The planning overhead matters. At rank 397, one adaptive analytic-ledger query is slower than one exact tilted-COS query. The crossover appears between ranks 397 and 797 in this prototype, while compiled reuse is much faster. No universal speedup claim is therefore made.



\section{Numerical verification}\label{sec:final-28}

All identities were compared against explicitly formed reference matrices.

\begin{table}[htbp]
\centering
\caption{Partition-invariance checks.}
\begin{tabular}{lrrr}
\toprule
partition & blocks & Gram relative error & action relative error\\
\midrule
coarse 4 & 4 & \(2.91\times10^{-16}\) & \(5.99\times10^{-16}\)\\
fine 64 & 32 & \(3.39\times10^{-16}\) & \(6.00\times10^{-16}\)\\
irregular & 9 & \(3.34\times10^{-16}\) & \(5.87\times10^{-16}\)\\
singleton-front & 21 & \(2.43\times10^{-16}\) & \(5.82\times10^{-16}\)\\
permuted, block 73 & 29 & \(3.52\times10^{-16}\) & \(5.89\times10^{-16}\)\\
\bottomrule
\end{tabular}
\end{table}

A hierarchical coarse-to-fine refinement gave relative discrepancy
\[
3.17\times10^{-16}.
\]

For an explicitly PSD polyphase operator
\[
H=L^\top L
\]
with block-bidiagonal \(L\), the smallest eigenvalue was approximately
\[
2.86,
\]
and the largest norm observed outside the exact \(k\)-step light cone was
\[
\boxed{0}.
\]

The QR-core update reproduced the full singular spectrum with relative
\(\ell_2\) discrepancy
\[
5.89\times10^{-16}.
\]



